\section{Rarefaction field on the characteristic surface}

On the characteristic surface $\mathcal C$ we have $Y=0$ and therefore
$U_-=U_+=\bar U(\tau,z)$.  The projective direction
\[
r(z)=\binom{z}{1}
\]
is an eigenvector of $DF(\bar U)$ associated with the characteristic speed
$s$:
\[
DF(\bar U)r=s r.
\]
A rarefaction curve $U(\xi)$ in state space is tangent to this selected
eigendirection, hence
\[
\frac{dU}{d\xi}=\alpha(\xi)r
\]
for some scalar reparametrization factor $\alpha$.

In the finite chart write
\[
\bar U(\tau,z)=U_0(z)+\tau T(z),
\qquad
T(z)=\partial_\tau\bar U=\binom{-Q(z)}{P(z)},
\]
and set $U_z=\partial_z\bar U$.  The implementation uses the
desingularized lifted field
\[
\boxed{
\dot\tau=\det(r,U_z),
\qquad
\dot z=\det(T,r)=-(Q+zP).
}
\]
Indeed, the two-dimensional determinant identity gives
\[
T\det(r,U_z)+U_z\det(T,r)=\det(T,U_z)r,
\]
so that
\[
\boxed{
D\bar U(\dot\tau,\dot z)=\det(T,U_z)r.
}
\]
Thus every regular integral curve of the lifted field projects to a
rarefaction curve in state space, up to reparametrization.  The factor
$\det(T,U_z)$ also makes explicit where the coordinate map itself becomes
singular; no division by this factor is used in the implementation.

For visualization only, the finite projective coordinate is compactified by
\[
\hat z=\frac{2}{\pi}\arctan z,
\qquad
\boxed{
\dot{\hat z}=\frac{2}{\pi(1+z^2)}\dot z.
}
\]
The compactification does not define the rarefaction field; it only changes
its displayed coordinate.

\subsection{Slow and fast characteristic sheets}

At a regular point of $\mathcal C$, the direction $r(z)$ selects the
eigenvalue $s$.  Since the system is $2\times2$, the other eigenvalue is
\[
\lambda_{\mathrm{other}}=\operatorname{tr}DF(\bar U)-s.
\]
Away from the coincidence locus $\mathcal E$, this orders the two sheets of
the characteristic surface:
\[
\boxed{
\mathcal C_s:\ s<\lambda_{\mathrm{other}},
\qquad
\mathcal C_f:\ s>\lambda_{\mathrm{other}}.
}
\]
On $\mathcal E$ the two speeds coincide and the slow/fast distinction
vanishes.  The function \texttt{characteristicSpectrum} implements exactly
this comparison; it does not infer the branch from the sign or magnitude of
$z$.

\subsection{Finite singularities}

A finite singularity of the lifted rarefaction field satisfies
\[
\boxed{\dot\tau=0,\qquad \dot z=0.}
\]
Because
\[
\dot z=-(Q+zP)
\]
depends only on $z$, the implementation first solves
\[
Q(z)+zP(z)=0.
\]
For the general quadratic flux this is the cubic equation
\[
\boxed{
b_{20}z^3+(b_{11}-a_{20})z^2+(b_{02}-a_{11})z-a_{02}=0.
}
\]
For every regular real root, $\dot\tau$ is affine in $\tau$, so the remaining
scalar equation is solved directly.  Classification is performed only after
the singular point has been found, using the Jacobian of
$(\dot\tau,\dot{\hat z})$ in the displayed finite chart.

The complementary chart $Z=1/z$ is required for singularities at projective
infinity.  Their local analysis and separatrices are implemented in the reciprocal chart below. They remain logically downstream of the definition of the rarefaction field: the chart analysis does not alter that definition.

\subsection{Singularities in the reciprocal chart}

The complementary chart is
\[
Z=\frac1z,
\qquad
\eta=\frac{\tau}{Z^2},
\]
and the lifted field used by the implementation is
\[
\dot\eta=U_Z^{(2)}-Z U_Z^{(1)},
\qquad
\dot Z=-Q(Z)Z-P(Z),
\]
where $P,Q$ are the reciprocal-chart coefficients and
$U_Z=\partial_Z\bar U(\eta,Z)$.
At projective infinity, $Z=0$, one obtains
\[
\boxed{\dot Z\big|_{Z=0}=-2b_{20}.}
\]
Consequently, a regular singularity of this chart can lie at projective
infinity only if $b_{20}=0$.  When $b_{20}=0$ and $b_{11}\neq0$, the chart
is regular at $Z=0$ and
\[
\dot\eta\big|_{Z=0}
=
(b_{11}-2a_{20})\eta
+
\frac{a_{10}b_{11}-a_{11}b_{10}-b_{01}b_{11}+2b_{02}b_{10}}
     {b_{11}^2}.
\]
If $2a_{20}\neq b_{11}$, the isolated singularity is therefore
\[
\boxed{
\eta_\infty=
\frac{a_{10}b_{11}-a_{11}b_{10}-b_{01}b_{11}+2b_{02}b_{10}}
     {b_{11}^2(2a_{20}-b_{11})},
\qquad Z=0.
}
\]
The Jacobian there is upper triangular.  Its diagonal entries, hence its
eigenvalues, are
\[
\boxed{
\lambda_1=b_{11}-2a_{20},
\qquad
\lambda_2=2(a_{20}-b_{11}).
}
\]
Thus the local type at infinity must be determined from this reciprocal
chart, not from the apparent behavior at the two ends of a compactified
finite-$z$ drawing.

For the legacy-wave preset used for regression tests,
\[
a_{20}=\frac92,\quad b_{11}=1,\quad b_{20}=0,\quad
b_{02}=-\frac1{10},\quad b_{10}=1,
\]
with the remaining relevant linear coefficients zero.  Hence
\[
\boxed{(\eta_\infty,Z)=\left(-\frac1{40},0\right)},
\qquad
\boxed{(\lambda_1,\lambda_2)=(-8,7)}.
\]
The projective singularity is therefore a \emph{saddle}.  A visualization
that makes this point look like a source or sink is not evidence of a
change of dynamical type; it indicates that the chart identification,
rendering, or separatrix continuation must be inspected.

\subsection{Local separatrices of the saddle at projective infinity}

Let $q_\infty=(\eta_\infty,0)$ be a hyperbolic saddle of the reciprocal-chart
field, with eigenvalues $\lambda_s<0<\lambda_u$ and corresponding eigenvectors
$v_s,v_u$.  The four local half-separatrices are initialized at
\[
q_\infty\pm\varepsilon v_s,
\qquad
q_\infty\pm\varepsilon v_u.
\]
To trace all four branches \emph{away} from the saddle, the unstable pair is
integrated in forward time and the stable pair in backward time.  Multiplying
the vector field by a positive normalization factor changes only its time
parametrization, not these integral curves.

For the regression preset, the stable eigendirection is tangent to the
projective boundary $Z=0$.  Consequently its two half-separatrices satisfy
\[
\boxed{Z\equiv0}
\]
locally (and, for this field, along their computed continuation).  They have
no representation in the finite $z$ chart because $z=1/Z$ is undefined on
that boundary.  In particular, mapping all points with $Z=0$ to
$\hat z=\pm1$ collapses the $\eta$ direction and creates the false visual
impression that a saddle separatrix is missing.

The unstable pair has $Z\neq0$ immediately away from the saddle, one branch
on each side of $Z=0$.  Wherever $Z\neq0$, its finite-chart coordinates are
recovered exactly by
\[
z=\frac1Z,
\qquad
\tau=\eta Z^2.
\]
Thus the reciprocal chart is not merely a numerical convenience: it is
necessary to display the complete local saddle geometry at projective
infinity.  The main finite-$z$ rendering must not invent a finite projection
for a branch contained in $Z=0$; a reciprocal/local chart is required to
visualize that branch.

\subsection{Projection of rarefaction curves to state space}

Let $\gamma(\sigma)=(\tau(\sigma),z(\sigma))\subset\mathcal C$ be an
integral curve of the lifted rarefaction field.  Since $Y=0$ on
$\mathcal C$, the canonical projections agree:
\[
 \pi_-(\gamma(\sigma))=\pi_+(\gamma(\sigma))
 =\bar U(\tau(\sigma),z(\sigma)).
\]
We therefore define the state-space rarefaction as the projection
\[
\boxed{\mathcal R(\sigma)=\bar U(\tau(\sigma),z(\sigma)).}
\]
No second differential equation is integrated in $(u,v)$.  From the lifted
field identity already established above,
\[
 \frac{d\mathcal R}{d\sigma}
 =\det(T,U_z)\,r(z),
\]
so the tangent of the projected curve is parallel to the selected
characteristic eigenvector.  Spectral classification is inherited pointwise:
segments in $\mathcal C_s$ project to $\mathcal R_s$ and segments in
$\mathcal C_f$ project to $\mathcal R_f$.  At $\mathcal E$ the two
eigenvalues coincide, so the slow/fast label is not assigned across that
point.

\subsection{A mesma órbita através da coincidência}

A classificação lenta ou rápida pertence aos \emph{trechos} da curva integral,
não necessariamente à órbita inteira.  Se uma órbita regular $\gamma$ cruza
transversalmente o locus de coincidência $\mathcal E$, existe um parâmetro
$\sigma_E$ tal que
\[
 p_E=\gamma(\sigma_E)\in\mathcal E,
\]
e, localmente, a mesma órbita pode ser escrita como
\[
\boxed{
\gamma=\gamma_s\cup\{p_E\}\cup\gamma_f
}
\]
(ou com a ordem $f,s$ invertida).  Assim, $\mathcal E$ não é uma condição de
parada do campo de rarefação: é o lugar onde a ordenação espectral que nomeia
$\mathcal C_s$ e $\mathcal C_f$ deixa de ser definida.

Numericamente, a órbita é integrada uma única vez. Quando o gap espectral
$s-\lambda_{\mathrm{other}}$ muda de sinal entre duas amostras consecutivas,
o evento é refinado por bisseção no segmento da órbita em coordenadas da
carta. O mesmo ponto refinado $p_E$ é colocado como extremidade dos dois
segmentos de visualização. Dessa forma não se introduz uma segunda órbita nem
uma descontinuidade artificial no espaço de estados.
